\section{Observer}
\label{sec:pat-observer}

\problemstatement{Define a one-to-many dependency so that when one object
(the subject) changes state, all its dependents (observers) are notified
and updated automatically -- the foundation of event systems and
publish/subscribe.}

\begin{patternbox}[When You Reach for It]
The trigger is \emph{``when X changes, several other things must react''} --
a spreadsheet cell and its charts, a stock price and its displays, a model
and its views. When one source has many interested listeners that come and
go, use Observer.
\end{patternbox}

\subsection*{Understanding the pattern}

The \textbf{Subject} maintains a list of \textbf{Observers} and offers
\texttt{subscribe}/\texttt{unsubscribe}. On a relevant change it calls
\texttt{notify()}, which invokes \texttt{update()} on every observer. The
subject knows only the observer \emph{interface}, so observers can be added
or removed at run time without the subject knowing their concrete types.
This is the classic Model-View decoupling.

\begin{ideabox}[Publish to a List of Subscribers]
Keep subscribers behind an interface and iterate them on change. Loose
coupling is the point: the subject broadcasts an event; who listens, and
how many, is not its concern and can vary dynamically.
\end{ideabox}

\subsection*{C++ implementation}

\begin{lstlisting}
#include <bits/stdc++.h>
using namespace std;

struct Observer {
    virtual void update(int price) = 0;
    virtual ~Observer() = default;
};

class Stock {                                   // subject
    vector<Observer*> observers;
    int price = 0;
public:
    void subscribe(Observer* o)   { observers.push_back(o); }
    void unsubscribe(Observer* o) {
        observers.erase(remove(observers.begin(), observers.end(), o),
                        observers.end());
    }
    void setPrice(int p) {
        price = p;
        for (Observer* o : observers) o->update(price);   // notify all
    }
};

struct Display : Observer {
    string name;
    Display(string n) : name(move(n)) {}
    void update(int price) override {
        cout << name << " sees price " << price << "\n";
    }
};

int main() {
    Stock stock;
    Display a("Phone"), b("Web");
    stock.subscribe(&a); stock.subscribe(&b);
    stock.setPrice(100);                        // both notified
    stock.unsubscribe(&a);
    stock.setPrice(120);                        // only Web notified
    return 0;
}
\end{lstlisting}

\begin{complexitybox}[Trade-offs]
\textbf{Pros.} Loose coupling between subject and observers; observers
added/removed at run time; supports broadcast/event architectures.
\textbf{Cons.} Notification order is unspecified; risk of update storms or
cycles; a dangling observer that forgot to unsubscribe causes bugs or
leaks (use weak references or careful lifetime management).
\end{complexitybox}

\begin{takeawaybox}
Observer wires a one-to-many ``notify on change'' relationship behind an
observer interface, the core of MVC and pub/sub. Interview keyword: ``when
this changes, update all those.'' The pitfalls to name are lapsed-listener
leaks and unspecified notification order.
\end{takeawaybox}
